\section{Introduction}
\begin{table}[ht]
\centering
\resizebox{\linewidth}{!}{
\begin{tabular}{lcccccccc}
\toprule
Paper & \#Claims & Claim Src & Natural? & Evidence Src & \multicolumn{2}{c}{No Leakage} & \(\text{Numerical}^{*}\) \\
\cmidrule(lr){6-7}
& & & & & Gold & Temporal & \\
\midrule
Thorne and Vlachos\cite{thorne-vlachos-2017-extensible} & 7k & KB & \xmark & KB & NA & NA & \xmark \\
Vlachos and Ridel\cite{vlachos-riedel-2015-identification} & 7k & KB & \xmark & KB & NA & NA & \cmark \\
FavIQ\cite{FAVIQ} & 188k & QA & \xmark & Wikipedia & NA & NA & \xmark \\
VitaminC\cite{VitaminC} & 326k & Wikipedia & \xmark & Wikipedia & NA & NA & \xmark \\
HOVER\cite{jiang2020hover} & 26k & QA & \xmark & Wikipedia & NA & NA & \xmark \\
SciFact\cite{scifact} & 1.4k & Science Articles & \xmark & Science Articles & NA & NA & \xmark \\
SciFactOpen\cite{wadden-etal-2022-scifact} & 279 & Science Articles & \xmark & Science Articles & NA & NA & \xmark \\
WICE\cite{WICE} & 2k & Wikipedia & \xmark & Wikipedia & NA & NA & \xmark \\
FEVER\cite{FEVER} & 185k & Wikipedia & \xmark & Wikipedia & NA & NA & \xmark \\
FEVEROUS\cite{FEVEROUS} & 87k & Wikipedia & \xmark & Wikipedia & NA & NA & \xmark \\
\bottomrule
MultiFC\cite{multifc} & 36k & FCW & \cmark & Open Web & \xmark & \xmark & \xmark \\
ClaimDecomp\cite{claimdecomp} & 1.2k & FCW & \cmark & Fact Check Articles & NA & NA & \xmark \\
FinFact\cite{finfact} & 3.4k & FCW & \cmark & Open Web & \xmark & \xmark & \xmark \\
WatClaimCheck\cite{watclaimcheck} & 34k & FCW & \cmark & Open Web & \cmark & \cmark & \xmark \\
LIAR\cite{liar} & 13k & FCW & \cmark & NA & NA & NA & \xmark \\
QABriefs\cite{qabriefs} & 8.8k & FCW & \cmark & Open web & \cmark & \xmark & \xmark \\
AviriTec\cite{averitec} & 4.5k & FCW & \cmark & Open Web & \cmark & \cmark & \xmark \\
Jandhagi and Pujara\cite{QVS} & 500 & News Articles & \cmark & Select Datasets & \cmark & \cmark & \cmark \\
COVIDFACT\cite{saakyan-etal-2021-covid} & 4.1k & Reddit & \cmark & Open Web & \cmark & \cmark & \xmark \\
QuanTemp\cite{quantemp} & 15k & FCW & \cmark & Open Web & \cmark & \xmark & \cmark \\
\rowcolor{blue!25}
\quantempplus{}(\textbf{ours}) & 15k & FCW & \cmark & Open Web & \cmark & \cmark & \cmark \\
\bottomrule
\end{tabular}
}
\caption{Comparison of \quantempplus{} to other related fact checking datasets. Leakage refers to either the presence of fact-checking articles (Gold) in the evidence corpus or the inclusion of evidence items published after the corresponding claim (Temporal). We also judge if a significant number of \(\text{Numerical}^{*}\) claims are present in the dataset. In the table, FCW refers to Fact-Checking Websites. }
\label{tab:dataset_all}
\vspace{-2em}
\end{table}

The spread of misinformation and disinformation has surged with the rise of advanced digital tools to disseminate information. Although manual verification of such claims are performed by experts, the sheer volume of misinformation motivates the need for automated fact-checking approaches\cite{cohen2011computation}. 

Of particular significance among the wide range of claims are numerical claims which comprise quantitative or temporal information \cite{quantemp}. These claims are significant because of the ``illusion of numeric truth effect" where claims with quantitative information provide a mirage of truth \cite{sagara2009consumer} potentially leading to catastrophic effects. For instance, unsubstantiated incomplete claims were made by Purdue Pharma to market their drug which is believed to be the kickstarter of the Opioid pandemic in America, resulting in the loss of over 500,000 lives to date\cite{van2009promotion}. Verification of numerical claims requires numerical reasoning ability which involves understanding numeric patterns and mathematical concepts like arithmetic, numerical estimation, and data interpretation.  While several works focus on automated verification of natural claims \cite{programFC,shah2022numerical,vlachos-riedel-2015-identification,shah2023concord,QVS}, only a handful focus on claims that require processing numerical information. However, these techniques only focus on the detection of numerical claims\cite{shah2023concord,shah2022numerical} or are restricted to specific simple statistical properties\cite{thorne-vlachos-2017-extensible,vlachos-riedel-2015-identification}. QuanTemp \cite{quantemp} was the first benchmark to focus on diverse numerical claims but is limited in evidence quality due to temporal leakage.

A typical human fact-checker utilizes various skills to verify a natural numerical claim. A given numerical claim can have different explicit and implicit information needs that first needs to be addressed.  Therefore, relevant evidence is gathered from various sources, including the open web. This evidence is then used to reason about and determine the veracity of a natural numerical claim.  The search process of a human fact-checker which involves decomposing the claim along different aspects is a crucial skill that forms the foundation of the verification process. Emulating a real-world setting is essential to aid in the development of automated methods that encompass such skills.

\begin{figure}[h!]
    \includegraphics[width=\textwidth]{figures/data_collection_ecir26-cropped.pdf}
    \caption{Data creation pipeline of \quantempplus{}}
    \label{fig:data-collection}
\end{figure}
\vspace{-1em}
Current datasets mainly feature synthetic claims with lexical biases \cite{FAVIQ,Chen2023ComplexCV}, while those with natural claims often suffer from gold or temporal leakage\cite{averitec} or lack focus on numerical claims (see Table 1.1). The presence of leaked fact verification articles in the corpus trivializes retrieval, while including evidence published after the claim makes the corpus unrealistic. Consequently, any automatic method designed under these conditions is likely to fail during real-time deployment. Hence, our goal is to provide a realistic dataset that addresses all these gaps, aiding the right development of automated methods for verifying natural numerical claims. 

Forming a large-scale dataset to address the above gaps through crowdsourcing can be challenging as obtaining manual annotations from expert fact-checkers is expensive\cite{expensive}. Therefore, we employ weak supervision to extend QuanTemp\cite{quantemp} and create
\textbf{\quantempplus{}}: a dataset consisting of about 15k natural numerical claims, an open domain corpus consisting of 165.7k records, and the corresponding evidence relevance and veracity labels of the claims. To improve the evidence quality collected from web, we propose \name{}, a claim decomposition that emulates the search process of human fact-checkers. Evidences collected using queries from \name{} helps in realistically evaluating and understanding retrieval and verification bottlenecks in automated fact-checking methods.

%\begin{table}[htb!!] % The placement specifier can be [htbp]
\centering
\scalebox{0.85}{
\begin{tcolorbox}[title= Example: Claim from \quantempplus{}] \textbf{Claim}:\texttt{Prime Minister Narendra Modi breached the \textcolor{blue}{election protocol} by addressing a rally in Howrah on \textcolor{orange}{April 6.}}


\textit{\textbf{Ideal Retrieval}}:
\begin{itemize}
    \item \ldots From cooch behar \ldots Howrah \ldots election rally on \textcolor{orange}{April 6}, pm Modi will address\ldots
    \item \ldots Election Silence Period\ldots India (\textcolor{blue}{48 to 24 hours} in advance of polling day and on polling day)\ldots 
    \item \ldots General elections were held in India in seven phases from \textcolor{brown}{11 April to 19 May 2019}\ldots
\end{itemize}
\textit{\textbf{Search Process}}:The given numerical claim requires multiple aspects to be addressed. Firstly, evidence proving that PM Modi did indeed rally on April 6 needs to be fetched. Then, the information on the election silence period and the polling schedule are to be retrieved. Given this information, the downstream verifier can then reason about the numerical aspects to form the verdict.


\end{tcolorbox}
}

\captionof{figure}{Example numerical claim from \quantempplus{}, requiring to recognize and fetch explicit and implicit evidence.}\label{intro:example} 
\end{table}%

We aim to address the following research questions:
\begin{enumerate}
\setlength\itemsep{1em}
\item \textbf{RQ1}: Does \name{} based decomposition help retrieve quality evidence from the web for the verification of natural numerical claims?
    \item \textbf{RQ2}: How do existing claim decomposition methods perform in terms of evidence retrieval to verify numerical claims in \quantempplus{}?
    \item \textbf{RQ3}: What is the downstream impact of these claim decomposition methods on the task of verification of numerical claims?

\end{enumerate}